% Part:sets-functions-relations
% Chapter: size-of-sets
% Section: enumerations

\documentclass[../../../include/open-logic-section]{subfiles}

\begin{document}

\olfileid[pt]{sfr}{siz}{enm}

\olsection{Enumerações e Conjuntos \usetoken{P}{enumerable}}

\begin{editorial}
  Esta seção discute enumerações de conjuntos, definindo-as como
  sobrejeções de $\PosInt$. Faz as coisas devagar, para leitores com
  pouca formação matemática. Uma versão alternativa, mais concisa,
  é dada na \olref[enm-alt]{sec}, que define enumerações
  de outro modo: como bijeções com $\Nat$ (ou um segmento inicial).
\end{editorial}

\begin{explain}
Já demos exemplos de conjuntos listando seus !!{element}s.
Discutamos agora, de maneira mais geral, como e quando podemos listar os
!!{element}s de um conjunto, mesmo que esse conjunto seja infinito.
\end{explain}

\begin{defn}[Enumeração, informalmente]
Informalmente, uma \emph{enumeração} de um conjunto~$A$ é uma lista (possivelmente
infinita) de !!{element}s de~$A$ tal que todo !!{element} de $A$
aparece na lista em alguma posição finita. Se $A$ tem uma
enumeração, então diz-se que $A$ é \emph{!!{enumerable}}.
\end{defn}

\begin{explain}
Alguns pontos sobre enumerações:
\begin{enumerate}
\item Contamos como enumerações apenas listas que têm um começo e
  em que todo !!{element} além do primeiro tem um único
  !!{element} imediatamente anterior. Em outras palavras, há
  apenas um número finito de elementos entre o primeiro !!{element} da
  lista e qualquer outro !!{element}. Em particular, isso significa que todo
  !!{element} de uma enumeração tem uma posição finita: o primeiro
  !!{element} tem posição~$1$, o segundo posição~$2$, etc.
\item Podemos ter enumerações diferentes do mesmo conjunto~$A$ que
  diferem pela ordem em que os !!{element}s aparecem: $4$, $1$,
  $25$, $16$,~$9$ enumera os (conjunto dos) cinco primeiros quadrados
  tão bem quanto $1$, $4$, $9$, $16$,~$25$.
\item Enumerações redundantes ainda são enumerações: $1$, $1$, $2$,
  $2$, $3$, $3$,~\dots{} enumera o mesmo conjunto que $1$, $2$,
  $3$,~\dots{}.
\item Ordem e redundância \emph{importam} quando especificamos uma
  enumeração: podemos enumerar os inteiros positivos começando com
  $1$, $2$, $3$, $1$, \dots{}, mas o padrão é mais fácil de ver quando
  enumerados da maneira padrão como $1$, $2$, $3$, $4$,~\dots
\item Enumerações devem ter um começo: \dots, $3$, $2$, $1$ não é
  uma enumeração dos inteiros positivos porque não tem primeiro
  !!{element}. Para ver como isso segue da definição informal,
  pergunte a si mesmo: ``em que posição na lista o número 76
  aparece?''
\item O seguinte não é uma enumeração dos inteiros positivos:
  $1$, $3$, $5$, \dots, $2$, $4$, $6$, \dots\@ O problema é que os
  números pares ocorrem nas posições $\infty + 1$, $\infty + 2$, $\infty +
  3$, em vez de em posições finitas.
\item O conjunto vazio é !!{enumerable}: é enumerado pela lista vazia!{}
\end{enumerate}
\end{explain}

\begin{prop}
  Se $A$ tem uma enumeração, tem uma enumeração sem
  repetições.
\end{prop}

\begin{proof}
  Suponha que $A$ tem uma enumeração $x_1$, $x_2$, \dots{} em que cada
  $x_i$ é um !!{element} de~$A$. Podemos remover repetições de uma
  enumeração removendo !!{element}s repetidos. Por exemplo, podemos
  transformar a enumeração em uma nova na qual listamos $x_i$ se
  ele é !!a{element} de~$A$ que não está entre $x_1$, \dots,
  $x_{i-1}$ ou removemos $x_i$ da lista se já aparece entre
  $x_1$, \dots,~$x_{i-1}$.
\end{proof}

O último argumento mostra que, para obter um bom controle sobre
enumerações e conjuntos !!{enumerable}s e provar coisas sobre eles,
precisamos de uma definição mais precisa. O seguinte a fornece.

\begin{defn}[Enumeração, formalmente]
Uma \emph{enumeração} de um conjunto $A \neq \emptyset$ é qualquer
função !!{surjective} $f \colon \PosInt \to A$.
\end{defn}

\begin{explain}
Convençamo-nos de que a definição formal e a definição informal
usando uma lista possivelmente infinita são equivalentes. Primeiro, qualquer
função !!{surjective} de $\PosInt$ para um conjunto~$A$ enumera~$A$.
Tal função determina uma enumeração conforme definido informalmente acima:
a lista $f(1)$, $f(2)$, $f(3)$, \dots. Como $f$ é !!{surjective},
todo !!{element} de~$A$ é garantidamente o valor de~$f(n)$ para
algum~$n \in \PosInt$. Logo, todo !!{element} de $A$ aparece em alguma
posição finita na lista. Como a função pode não ser
!!{injective}, a lista pode ser redundante, mas isso é aceitável (como
observado acima).

Por outro lado, dada uma lista que enumera todos os !!{element}s
de~$A$, podemos definir uma função !!{surjective} $f\colon \PosInt \to A$
fazendo $f(n)$ ser o $n$-ésimo !!{element} da lista, ou o último
!!{element} da lista se não há $n$-ésimo !!{element}. O único
caso em que isso não produz uma função !!{surjective} é quando $A$ é
vazio, e portanto a lista é vazia. Assim, toda lista não vazia
determina uma função !!{surjective} $f\colon \PosInt \to A$.
\end{explain}

\begin{defn}
  \ollabel{defn:enumerable}
  Um conjunto~$A$ é !!{enumerable} sse é vazio ou tem uma enumeração.
\end{defn}

\begin{ex}
Uma função que enumera os inteiros positivos ($\PosInt$) é simplesmente a
função identidade dada por $f(n) = n$. Uma função que enumera os
números naturais $\Nat$ é a função $g(n) = n - 1$.
\end{ex}

\begin{ex}
As funções $f\colon \PosInt \to \PosInt$ e $g \colon \PosInt \to
\PosInt$ dadas por
\begin{align*}
f(n) & = 2n \text{ e}\\
g(n) & = 2n - 1
\end{align*}
enumeram os inteiros positivos pares e os inteiros positivos ímpares,
respectivamente. Contudo, nenhuma das funções é uma enumeração de
$\PosInt$, já que nenhuma é !!{surjective}.
\end{ex}

\begin{prob}
Defina uma enumeração dos quadrados positivos $1$, $4$, $9$, $16$, \dots
\end{prob}

\begin{ex}
A função $f(n) = (-1)^{n} \lceil \frac{(n-1)}{2}\rceil$ (onde
$\lceil x \rceil$ denota a função \emph{teto}, que arredonda
$x$ para cima até o inteiro mais próximo) enumera o conjunto dos
inteiros~$\Int$. Note como $f$ gera os valores de $\Int$
``saltando'' de um lado para outro entre inteiros positivos e negativos:
\[
\begin{array}{c c c c c c c c}
f(1) & f(2) & f(3) & f(4) & f(5) & f(6) & f(7) & \dots \\ \\
- \lceil \tfrac{0}{2} \rceil & \lceil \tfrac{1}{2}\rceil & - \lceil \tfrac{2}{2} \rceil & \lceil \tfrac{3}{2} \rceil & - \lceil \tfrac{4}{2} \rceil  & \lceil \tfrac{5}{2}
\rceil & - \lceil \tfrac{6}{2} \rceil & \dots \\ \\
0 & 1 & -1 & 2 & -2 & 3 & \dots
\end{array}
\]
Você também pode pensar em $f$ como definida por casos da seguinte forma:
\[
f(n) = \begin{cases}
  0 & \text{se $n = 1$}\\
  n/2 & \text{se $n$ é par}\\
  -(n-1)/2 & \text{se $n$ é ímpar e $>1$}
  \end{cases}
\]
\end{ex}

\begin{prob}
  Mostre que se $A$ e $B$ são !!{enumerable}s, então $A \cup B$ também é. Para
  isso, suponha que há funções !!{surjective}s $f\colon \PosInt \to
  A$ e $g\colon \PosInt \to B$, e defina uma função
  !!{surjective}~$h\colon \PosInt \to A \cup B$ e mostre que ela é
  !!{surjective}. Considere também os casos em que $A$ ou~$B = \emptyset$.
\end{prob}

\begin{prob}
  Mostre que se $B \subseteq A$ e $A$ é !!{enumerable}, então $B$ também é. Para
  isso, suponha que há uma função !!{surjective} $f\colon \PosInt \to
  A$. Defina uma função !!{surjective}~$g\colon \PosInt \to B$ e mostre
  que ela é !!{surjective}. O que acontece se $B = \emptyset$?
\end{prob}

\begin{prob}
  Mostre por indução em $n$ que se $A_1$, $A_2$, \dots, $A_n$ são todos
  !!{enumerable}s, então $A_1 \cup \dots \cup A_n$ também é. Você pode assumir o
  fato de que se dois conjuntos $A$ e~$B$ são !!{enumerable}s, então $A \cup
  B$ também é.
\end{prob}

Embora seja talvez mais natural ao listar os !!{element}s de um
conjunto começar a contar a partir do $1$º !!{element}, os matemáticos gostam
de usar os números naturais~$\Nat$ para contar coisas. Eles
falam do $0$º, $1$º, $2$º, e assim por diante, !!{element} de uma lista.
Correspondentemente, podemos definir uma enumeração como uma função
!!{surjective} de $\Nat$ para~$A$. É claro que as duas definições são
equivalentes.

\begin{prop}\ollabel{prop:enum-shift}
  Há !!a{surjection} $f\colon \PosInt \to A$ sse há
  !!a{surjection} $g\colon \Nat \to A$.
\end{prop}

\begin{proof}
  Dada !!a{surjection} $f\colon \PosInt \to A$, podemos definir $g(n) =
  f(n+1)$ para todo $n \in \Nat$. É fácil ver que $g\colon \Nat
  \to A$ é !!{surjective}. Reciprocamente, dada !!a{surjection} $g\colon
  \Nat \to A$, defina $f(n) = g(n-1)$.
\end{proof}

Isso nos dá o seguinte resultado:

\begin{cor}\ollabel{cor:enum-nat}
Um conjunto $A$ é !!{enumerable} sse é vazio ou há
uma função !!{surjective} $f\colon \Nat \to A$.
\end{cor}

Discutimos acima que uma lista de !!{element}s de um conjunto~$A$ pode ser
transformada em uma lista sem repetições. Isso também vale para
enumerações, mas é um pouco mais difícil de formular e provar rigorosamente. Qualquer
função $f\colon \PosInt \to A$ deve estar definida para todo $n \in
\PosInt$. Se há apenas um número finito de !!{element}s em~$A$, então claramente não podemos ter uma função definida nos infinitos
!!{element}s de~$\PosInt$ que toma como valores todos os !!{element}s
de~$A$ mas nunca toma o mesmo valor duas vezes. Nesse caso, isto é, no
caso em que a lista sem repetições é finita, devemos escolher
um domínio diferente para~$f$, um com apenas um número finito de !!{element}s.
Não ter repetições significa que $f$ deve ser !!{injective}. Como também é
!!{surjective}, estamos procurando !!a{bijection} entre algum
conjunto finito $\{1, \dots, n\}$ ou $\PosInt$ e~$A$.

\begin{prop}\ollabel{prop:enum-bij}
Se $f\colon \PosInt \to A$ é !!{surjective} (isto é, uma enumeração
de~$A$), há !!a{bijection} $g\colon Z \to A$ onde $Z$ é
$\PosInt$ ou $\{1, \dots, n\}$ para algum~$n \in \PosInt$.
\end{prop}

\begin{proof}
  Definimos a função $g$ recursivamente: seja $g(1) = f(1)$. Se $g(i)$
  já foi definido, seja $g(i+1)$ o primeiro valor de $f(1)$,
  $f(2)$, \dots{} que ainda não está entre $g(1)$, \dots, $g(i)$, se houver
  um. Se $A$ tem apenas $n$ !!{element}s, então $g(1)$, \dots, $g(n)$ estão todos
  definidos, e assim definimos uma função $g\colon \{1, \dots, n\}
  \to A$. Se $A$ tem infinitos !!{element}s, então para qualquer $i$
  deve haver !!a{element} de~$A$ na enumeração $f(1)$, $f(2)$,
  \dots, que ainda não está entre $g(1)$, \dots, $g(i)$. Nesse
  caso definimos uma função $g\colon \PosInt \to A$.

  A função $g$ é !!{surjective}, já que qualquer elemento de~$A$ está
  entre $f(1)$, $f(2)$, \dots{} (pois $f$ é !!{surjective}) e assim
  eventualmente será um valor de~$g(i)$ para algum~$i$. Também é
  !!{injective}, já que se houvesse $j < i$ tal que $g(j) = g(i)$,
  então $g(i)$ já estaria entre $g(1)$, \dots, $g(i-1)$, contrariamente
  à forma como definimos~$g$.
\end{proof}

\begin{cor}\ollabel{cor:enum-nat-bij}
Um conjunto $A$ é !!{enumerable} sse é vazio ou há !!a{bijection}
$f\colon N \to A$ onde $N = \Nat$ ou $N = \{0, \dots, n\}$ para
algum $n \in \Nat$.
\end{cor}

\begin{proof}
$A$ é !!{enumerable} sse $A$ é vazio ou há !!a{surjective}
$f\colon \PosInt \to A$. Pela \olref{prop:enum-bij}, o último vale
sse há uma função !!{bijective}~$f\colon Z \to A$ onde $Z =
\PosInt$ ou $Z = \{1, \dots, n\}$ para algum $n \in \PosInt$. Pelo
mesmo argumento da prova da \olref{prop:enum-shift}, isso por sua vez
é o caso sse há !!a{bijection} $g\colon N \to A$ onde
$N = \Nat$ ou $N = \{0, \dots, n-1\}$.
\end{proof}

\begin{prob}
  De acordo com a \olref[sfr][siz][enm]{defn:enumerable}, um conjunto $A$ é
  !!{enumerable} sse $A = \emptyset$ ou há !!a{surjective} $f\colon
  \PosInt \to A$. Também é possível definir ``conjunto !!{enumerable}''
  precisamente por: um conjunto é !!{enumerable} sse há uma função
  !!{injective} $g\colon A \to \PosInt$. Mostre que as definições são
  equivalentes, isto é, mostre que há uma função !!{injective}
  $g\colon A \to \PosInt$ sse $A = \emptyset$ ou há
  !!a{surjective} $f\colon \PosInt \to A$.
\end{prob}

\end{document}
